\section{Temporal validity conflicts and forgetting}\label{sec:C14:validity}
Prospective memory requires more than answering a question about the past: a future event must activate a previously established constraint. TriggerBench separates an early constraint from a later trigger using intervening dialogue, matched retrospective questions, overloaded positive triggers, and negative controls~\cite{zhang2026triggerbench}. Its distinction motivates an implementation object $(g,u,b)$: $g$ is a trigger predicate over observed events, $u$ is the intended intervention, and $b$ is an applicability boundary including expiry and authority. At an event, evaluate the guard and current boundary before proposing the intervention. This is a proposed representation, not a claim that the benchmark implements an event engine. Evaluate missed triggers and false activations separately; unconditional reminders can improve recall while failing negative cases. A reminder also cannot confer authority to execute the underlying action.

Memory usefulness depends on whether an assertion remains valid when it is used. A correct historical fact can be a wrong current instruction. A past failed attempt can be useful as an episode while its error cause has been fixed. Temporal validity should therefore be represented at the proposition or procedure level, not inferred from how recently a record was retrieved. Retrieval frequency measures access, not truth.

Use valid-time intervals $[t_b,t_e)$ where meaningful, and preserve observation time separately. If a user states at time $t_2$ that a preference changed at $t_1<t_2$, the current resolver should incorporate the earlier effective change without pretending the agent knew it before $t_2$. This supports both present decisions and historical audit. A last-write-wins store ordered by ingestion can produce the wrong historical answer, while a store ordered only by effective time can lose what was actually known at a decision.

Conflicts can be factual, temporal, procedural, or authoritative. Factual conflict concerns incompatible descriptions of the same state. Temporal conflict can disappear once intervals are distinguished. Procedural conflict concerns strategies with different preconditions. Authority conflict concerns who may direct an action. A generic contradiction prompt can help identify conflicts, but the resolver's rule must reflect the class. Majority voting over repeated model-generated statements is especially weak evidence of authority or truth.

Forgetting has several meanings: preventing retrieval, removing a record, deleting derivatives, expiring an active preference, or modifying parameters. An external memory system should state which operation it implements. Removing a row from one index does not establish erasure from raw logs, summaries, backups, or downstream services. Conversely, revoking future use can be a meaningful operational guarantee even when some audit record remains. Claims should follow the actual deletion and retention contract.

A proposed invalidation protocol marks the source assertion obsolete, finds derived summaries and skills that depend on it, invalidates their reusable status, and rebuilds any current view that included them. This assumes a recorded dependency graph. Missing provenance limits completeness. A safe fallback can quarantine a larger region of memory at the cost of reduced utility. That tradeoff should be measured rather than hidden behind a claim of perfect forgetting.

Longitudinal evaluation should include corrected facts, changed preferences, expired permissions, and requests to forget. Score not only the desired new answer but also whether the agent stops using the old record in unrelated later tasks. Include negative controls where an update should not affect another project or user. These tests assess selective revision rather than indiscriminate amnesia. Their relevance increases as sessions, teams, and model versions share a growing store.

